% LSPV-1 REPORT::PROMFAST1::2026-10-12T12:16:33Z
% !TEX program = xelatex
\documentclass[11pt]{article}
\usepackage[margin=0.88in]{geometry}
\usepackage{amsmath,amssymb,booktabs,array,enumitem,xurl,hyperref}
\hypersetup{colorlinks=true,linkcolor=blue,urlcolor=blue}
\newcommand{\Pclass}{\mathsf P}
\newcommand{\NPclass}{\mathsf{NP}}
\newcommand{\PromiseP}{\mathsf{PromiseP}}
\newcommand{\OPEN}{\textbf{OPEN}}
\newcommand{\EXACT}{\textbf{PROVED/EXACT}}
\newcommand{\KNOWN}{\textbf{KNOWN THEOREM}}
\newcommand{\COMP}{\textbf{COMPUTATIONAL}}
\newcommand{\RETIRED}{\textbf{FALSIFIED/RETIRED}}
\title{PROMFAST1: Promise fast paths exactly classify recognizable candidate cores}
\author{LSPV-1 / UNIVERSAL\_CORE\_LOWER\_BOUND}
\date{12 October 2026, 12:16 UTC}

\begin{document}
\maketitle

\section*{Executive result}
\textbf{P versus NP remains OPEN.}  This run proves an exact boundary for a
frequent proposed universal-core obligation.  Let $L\subseteq\{0,1\}^*$ be
recursive and let $C\subseteq\{0,1\}^*$ be polynomial-time decidable.  There
exists a total exact decider for $L$ that runs in polynomial time on every
input in $C$ if and only if the promise restriction of $L$ to $C$ is in
$\PromiseP$.  Therefore any polynomial-time recognizable all-YES or all-NO
subfamily admits a constant-output fast path inside an exact global decider
and cannot be a universal time-hard core.

This result does not show that the active signed Monotone $3$SAT-$(2,2)$
domain is easy.  That domain is a mixed NP-complete promise.  Nor does the
result contradict Lynch's complexity-core theorem: a Lynch core is merely
recursive, need not be polynomial-time recognizable, and need not lie inside
the yes-instances.  The new theorem instead closes the tempting repair
``choose an explicit recognizable proper yes-core and demand that every exact
decider be slow there.''

\section{Claim card: exact model and quantifiers}
Fix a deterministic multitape Turing-machine model, a self-delimiting binary
encoding, and encoded length $N=|x|$.  Let
$\operatorname{Dec}(L)$ denote the total deterministic uniform exact
deciders for recursive $L$.  All parsing, promise recognition, output, and
fallback work is charged.

For a fixed polynomial-time decider $R_C$ of $C$, say that $A$ is
\emph{polynomial on $C$} when
\[
 \exists p\in\mathbb N[X]\;\forall x\in C:\quad T_A(x)\le p(|x|).
\]
Define $L\upharpoonright C\in\PromiseP$ to mean that a total deterministic
polynomial-time machine $B$ agrees with $L$ on every $x\in C$; its
off-promise output is unrestricted.  Requiring $B$ to be total and globally
polynomial removes any convention about off-promise halting.

\paragraph{Target tested.}
Can a fixed polynomial-time recognizable clean family force every exact
global decider to expose or pay superpolynomial time for compatibility?

\paragraph{Weakest useful theorem.}
Classify that all-decider time obligation before attempting an internal-state
or communication lower bound.  A single promise-easy counterexample decider
falsifies the universal claim.

\paragraph{SAME\_REALIZATION.}
The relevant realization is the complete wrapper machine constructed below.
It recognizes $C$, takes either the promise fast path or the exact fallback,
and outputs one bit.  There is no free extraction, uncharged oracle, or mixing
of the fast path with a different execution.

\section{PROMISE-FASTPATH-EQUIVALENCE1}
\paragraph{Theorem.}
For every recursive $L$ and every polynomial-time decidable $C$,
\[
\boxed{
 \bigl(\exists A\in\operatorname{Dec}(L)\;\exists p\;
       \forall x\in C:T_A(x)\le p(|x|)\bigr)
 \iff L\upharpoonright C\in\PromiseP .}
\]

\paragraph{Proof.}
For the forward implication, let $A$ be exact and polynomial on $C$.
Machine $B$ first runs $R_C$.  If $x\notin C$, it outputs $0$; if
$x\in C$, it runs $A(x)$.  The recognizer is globally polynomial, the
second branch is polynomial by hypothesis, and $B$ is correct on the
promise.

For the reverse implication, let $B$ be a total polynomial-time promise
decider and let $D$ be any total recursive decider for $L$.  Define
\[
 A(x)=\begin{cases}
 B(x),&R_C(x)=1,\\
 D(x),&R_C(x)=0.
 \end{cases}
\]
Then $A$ is total and exact on all strings, while its execution on $C$
uses only $R_C$ and $B$, hence is polynomial. \hfill\(\square\)

If $C$ contains arbitrarily long strings, absorbing finitely many exceptional
inputs into a polynomial constant gives the equivalent cofinal negation
\[
 \forall A\in\operatorname{Dec}(L)\;\forall p\;\forall N_0\;
 \exists x\in C:\quad |x|\ge N_0\ \wedge\ T_A(x)>p(|x|)
 \iff L\upharpoonright C\notin\PromiseP .
\]
This is \EXACT, pending independent review.  It is a classification theorem,
not a proof that either side holds for an NP-hard promise.

\paragraph{Corollary (recognizable proper cores fail).}
If $C\subseteq L$ or $C\subseteq\overline L$ and $C\in\Pclass$, then
the promise problem is solved by a constant output.  The wrapper above is an
exact global decider polynomial on $C$.  Thus a polynomial-time recognizable
monochromatic family cannot be a Lynch polynomial complexity core and cannot
witness an all-decider time lower bound.  This corollary is \EXACT, pending
independent review.

\section{Relation to the explicit clean core}
The active syntactic domain $K=\mathcal H^{\pm}_{2,2}$ consists of
duplicate-free proper signed Monotone $3$SAT-$(2,2)$ encodings, with a
fixed malformed-input convention.  D\"ocker's theorem makes the restricted
satisfiability problem NP-complete.  Consequently $K$ is polynomial-time
recognizable but is not monochromatic, and
\[
 \mathsf{SAT}\upharpoonright K\in\PromiseP
 \iff \mathsf{SAT}\in\Pclass
 \iff \Pclass=\NPclass .
\]
The first equivalence uses the standard reduction into the NP-complete
restricted problem and the trivial restriction of a global SAT algorithm.
Hence the universal cofinal time target on this particular core remains
separation-strength and \OPEN.  PROMFAST1 neither proves nor falsifies it.

The distinction matters: a \emph{clean promise domain} can be explicit and
mixed, whereas a \emph{proper yes-instance core} is monochromatic.  Calling
both objects a ``core'' had permitted an invalid transfer between them.

\section{Finite hostile-control audit}
The deterministic script
\path{experiments/review/promise_fastpath_core_audit.py} exhaustively tested
all $256$ Boolean languages and all $256$ promise subsets of an
eight-element universe.  For each of $65{,}536$ language/promise pairs it
constructed the promise-fast-path/fallback wrapper and checked global exactness
and fast-path agreement.  It also registered 2SAT, Horn, affine/Tseitin,
bounded-treewidth, parity, signed-clean-core, all-YES, all-NO, padding,
redundant-clause, independent-product, alternate-encoding, one-bit-output,
DIRECT, COPY, DEAD, DELAY, INTERP, and RECOMP controls.

The run completed in $0.339$ seconds without chunking and performed
$786{,}451$ checks with zero failures.  Script SHA-256 is
\texttt{82219cee244ab60d8b87ccbbfa7af58a7232a36fdc1745e55fb6aab5326f77bd};
result SHA-256 is
\texttt{847a0d0d1bfc7c643e5dd2c1a31b0d82370e36d74a70dd942d71a4dab81247c6}.
The result is at
\path{results/raw/review/2026-10-12_promise_fastpath_core.json}.
This is \COMP{} only.  Every finite promise has a finite table, so the audit
tests construction bookkeeping and does not extrapolate asymptotics.

\section{Hostile controls, wrappers, and one-bit output}
\begin{itemize}[leftmargin=*]
\item On any recognizable Horn, 2SAT, affine/Tseitin, or fixed-treewidth
subfamily with a polynomial promise solver, the theorem builds an exact global
decider using that semantic shortcut.  No compatibility table is exposed.
\item All-YES and all-NO promises use a constant output.  This is the sharpest
counterexample to a recognizable proper-core obligation.
\item Padding, redundant clauses, independent products, and alternate
encodings do not create hardness when both the transformed promise and its
solver remain polynomial-time recognizable/computable; recognition and
translation costs are charged.
\item COPY, DEAD, DELAY, INTERP, and RECOMP can alter traces and internal
objects.  The existence theorem needs only one fast exact wrapper, so an
all-decider claim is already false when that wrapper exists.
\item The output is one decision bit.  Witness reconstruction, compatibility
extraction, communication, and memory exposure are absent unless a separate
model forces them.
\end{itemize}

\section{Primary-literature and barrier-transfer preflight}
The two major obstacles were (i) distinguishing a promise-easy recognizable
subfamily from a Lynch complexity core, and (ii) transferring semantic
structure to a charged resource of every arbitrary computation.
\begin{itemize}[leftmargin=*]
\item Even--Selman--Yacobi's promise-problem formalism is a
\textbf{directly transferable definition}: correctness is required on the
declared promise.  PROMFAST1 strengthens the machine convention to total
global polynomial time, which the recognizer wrapper supplies.
\item Lynch's polynomial-complexity-core theorem is a \textbf{directly
transferable theorem} for recursive languages outside P.  It does not say
that the core is polynomial-time recognizable or monochromatic.  PROMFAST1
therefore complements rather than reproves MODELBRIDGE-0024.
\item D\"ocker's NP-completeness theorem is directly transferable to the
active signed clean promise after fixing the encoding.
\item FLMV time--space lower bounds, branching-program and proof-system
bounds, VLSI congestion, communication cuts, control observability, and
local-to-global gluing are \textbf{useful structural analogies requiring a
new transfer theorem}.  They constrain declared models, not arbitrary
one-bit deciders.
\item Treating promise hardness, proof size, state count, or representation
size as ordinary decision time without a model-preserving compiler is a
\textbf{misleading analogy / failed transfer}.
\end{itemize}

The proof is elementary and relativizes and algebrizes.  It constructs no
large constructive circuit property and neither invokes nor escapes Natural
Proofs.  A positive P-versus-NP result still needs a nonrelativizing and
nonalgebrizing hinge or another barrier-crossing argument.

\section{Post-quantum technique ledger}
\begin{itemize}[leftmargin=*]
\item Schur squares, puncturing/shortening/duals, quadratic-relation ranks,
Hilbert functions, syzygies, and Betti numbers: \textbf{FAILED TRANSFER} to
arbitrary one-bit computations; hidden algebraic structure may instead make
the promise easy.
\item Collision entropy for a declared map $M:A\to F$, including
$|\operatorname{im}M|\ge |A|^2/\operatorname{Coll}(M)$:
\textbf{EXACT OBJECT/THEOREM TRANSFER} only for that fixed map.  PROMFAST1
shows no such map is forced by the promise interface.
\item LWE-style external extraction and worst-case/average-case architecture:
\textbf{ANALOGY/CANDIDATE needing a new bridge}; known NP barriers and all
query/reconstruction costs must be audited.
\item Special-soundness/collision extraction:
\textbf{ANALOGY/CANDIDATE} only with a construction-specific theorem that two
conflicting openings yield an unconditional forbidden witness.
\item Cryptographic binding, hashes, and random-oracle methods:
\textbf{FAILED TRANSFER} for an unconditional proof unless every assumption
is removed.
\end{itemize}

\section{CORPUS REVIEWED}
The current GitHub tree was enumerated.  This run read root
\path{main.tex}, actual \path{main4.tex} and \path{main5.tex},
\path{AGENTS.md}, \path{CODEX.md}, the complete pnp-research skill and its
protocol/mode references, and all current task specifications.  It read the
newest verified MAIN SYNTHESIS-0023; latest LSPV-1 CIRCEQUIV1, LSPV-2
REPDIRSUM1, LSPV-3 ANTISHARE1, and LSPV-4 MODELBRIDGE-0024; and materially
relevant COREEQUIV1, RESOURCEAXIS1, ALL-EASY-ISLAND-RESIDUAL-COLLAPSE1, and
complexity-core history surfaced by repository search.  The underlying
HEC, VSC-HEC, G-PPC-HEC, CRS-PPC, RCCN, and CCSPC causal components, current
governance, and cited experiment/result surfaces were inspected rather than
accepted only through peer summaries.

The immediately preceding LSPV-5 cycle had no newer verified canonical MAIN
after SYNTHESIS-0023 despite multiple scheduled opportunities.  Repository
history and successful specialist writes ruled out a repository/path-wide
permission failure, but no exact missing LSPV-5 scientific body was available.
A nonblocking recovery request preserving the no-reconstruction rule was
written before substantive work.  The failure remains classified UNKNOWN
with missing-body evidence; a governance escalation records the smallest
technical recovery procedure.

\section{DUPLICATION / CONTRADICTION AUDIT}
COREEQUIV1 already classified the all-decider worst-case time target on the
NP-hard signed clean domain as equivalent to $P\ne NP$.  MODELBRIDGE-0024
used Lynch to classify existence of one recursive almost-everywhere hard set.
SYNTHESIS-0021 showed that subtracting all polynomially recognizable
monochromatic islands is vacuous because singletons (and padded versions)
cover the domain.  PROMFAST1's nonduplicate increment is the exact
\emph{one-fixed-promise} equivalence: a polynomial-time recognizable promise
supports a fast exact-global wrapper precisely when its restricted decision
problem is in Promise-P.  It also gives the explicit monochromatic-core
counterexample.  No verified theorem is contradicted.

\section{WHAT CHANGED RELATIVE TO ROOT MAIN}
Root MAIN already states that the signed clean-core time target is
separation-strength and that internal extraction is not forced.  This report
adds the exact promise-fast-path boundary, retires polynomial-time
recognizable proper yes/no cores as universal time-hard candidates, and
clarifies that the active mixed NP-hard promise survives.  Actual
\path{main4.tex} and \path{main5.tex} remain partial ARC/WRAPPER construction
manuscripts and are not silently redefined as arbitrary-decider compilers.

\section{ADOPTED / REJECTED / RETIRED PRIOR COMPONENTS}
\begin{itemize}[leftmargin=*]
\item \textbf{ADOPTED:} corrected signed NP-complete promise domain; encoded
length; full charging; exact global fallback; SAME\_REALIZATION; hostile
wrappers; Lynch's recursive-core boundary.
\item \textbf{REJECTED:} semantic-state count as runtime; named extraction;
free promise recognition; nonuniform per-length tables; padding dilution;
transfer from restricted proof/communication models without a compiler.
\item \textbf{FALSIFIED/RETIRED:}
\texttt{P-RECOGNIZABLE-PROPER-CORE1}: a polynomial-time recognizable
monochromatic family cannot be time-hard for every exact global decider.
\item \textbf{OPEN:} promise hardness of the active mixed NP-complete core;
a wrapper-invariant full-language transfer; and a model-matched unconditional
landing for the same ordinary realization.
\end{itemize}

\section{STALE OR OVERSTATED CLAIMS FOUND}
Any use of ``proper core'' to mean a polynomial-time recognizable subset of
yes-instances while simultaneously claiming all-decider time hardness is
false by the constant-output wrapper.  Any claim that Lynch provides such a
proper recognizable core is overstated.  Conversely, PROMFAST1 must not be
overstated to make every recognizable mixed promise easy: on the active
NP-hard promise, its right-hand side is exactly the unresolved separation.

\section{Strict status ledger and handoff}
\begin{center}
\begin{tabular}{p{0.48\linewidth}p{0.20\linewidth}p{0.24\linewidth}}
\toprule
Item & Status & Scope \\
\midrule
Promise-problem formalism & \KNOWN & Even--Selman--Yacobi\\
Signed clean-core NP-completeness & \KNOWN & D\"ocker\\
Lynch complexity-core theorem & \KNOWN & Recursive languages outside P\\
\texttt{PROMISE-FASTPATH-EQUIVALENCE1} & \EXACT & Review pending\\
Recognizable monochromatic-core impossibility & \EXACT & Review pending\\
Finite wrapper audit & \COMP & 786,451 checks\\
\texttt{P-RECOGNIZABLE-PROPER-CORE1} & \RETIRED & Constant fast path\\
Active mixed clean-core time lower bound & \OPEN & Equivalent to $P\ne NP$\\
P versus NP & \OPEN & No unconditional separation\\
\bottomrule
\end{tabular}
\end{center}

\paragraph{Precise surviving target.}
Do not seek a polynomial-time recognizable proper yes-core.  Either prove the
mixed signed clean promise outside Promise-P---which is separation-strength---
or construct an external, wrapper-invariant, full-language relation from every
hypothetical polynomial SAT decider and land the same fully charged ordinary
realization in one cited unconditional forbidden region.  Arbitrary polynomial
exponents, alternate encodings, semantic shortcuts, and all hostile wrappers
must survive.  This target remains \OPEN.

\begin{thebibliography}{9}
\bibitem{ESY84} S. Even, A. L. Selman, and Y. Yacobi,
``The Complexity of Promise Problems with Applications to Public-Key
Cryptography,'' \emph{Information and Control} 61(2), 159--173, 1984.
\bibitem{Lynch75} N. A. Lynch, ``On Reducibility to Complex or Sparse Sets,''
\emph{Journal of the ACM} 22(3), 341--345, 1975.
\bibitem{Docker20} J. D\"ocker, ``Monotone 3-Sat-(2,2) is NP-complete,''
\emph{Theoretical Computer Science} 802, 1--9, 2020.
\bibitem{BGS75} T. Baker, J. Gill, and R. Solovay,
``Relativizations of the P=?NP Question,'' \emph{SIAM J. Comput.} 4(4),
431--442, 1975.
\bibitem{AW09} S. Aaronson and A. Wigderson, ``Algebrization: A New Barrier
in Complexity Theory,'' \emph{ACM TOCT} 1(1), 2009.
\bibitem{FLMV05} L. Fortnow, R. Lipton, D. van Melkebeek, and A. Viglas,
``Time-space lower bounds for satisfiability,'' \emph{J. ACM} 52(6),
835--865, 2005.
\end{thebibliography}

\noindent\textbf{Final verification sentinel:} report body complete; exact
quantifiers, encoded length, promise recognition, fallback, and all costs
stated; corpus and predecessor audited; strict statuses preserved; P versus
NP remains \OPEN.
% END_OF_VERIFIED_REPORT::LSPV-1::PROMFAST1
\end{document}
